\subsection{Unified Depth Fusion Frameworks}

Current systems are specialized for a single depth modality.  A unified
framework that can adaptively fuse whatever depth sources are
available---LiDAR, iToF, monocular depth, or combinations thereof---would
maximize flexibility across hardware platforms.  The depth completion
network of Gaussian-LIC2, which fuses RGB appearance with sparse LiDAR,
hints at this direction but remains LiDAR-specific.

\subsection{Learned Front-Ends for Gaussian SLAM}

Most Gaussian SLAM systems use classical odometry front-ends (IESKF,
DROID-SLAM, ORB-SLAM3).  Recent work on streaming multi-camera odometry
with frozen 3D foundation models (StreamRig~\cite{wei2026_streamrig}) and
learned visual-inertial dense BA (VIGS-SLAM's custom CUDA kernels)
suggests that learned front-ends can provide more robust correspondences
for Gaussian initialization, especially under challenging conditions.

\subsection{Memory-Bounded Gaussian Mapping}

Large-scale outdoor mapping with 3DGS generates millions of Gaussians,
exceeding GPU memory.  Current solutions include sliding-window
optimization (GS-LIVO), submap management (RMGS-SLAM, LV-GS SLAM), and
budgeted LOD (Budgeted-GS~\cite{wang2026_budgetedgs}).  A principled
memory-bounded Gaussian allocation strategy---driven by uncertainty
(BayesianGS-SLAM~\cite{kang2026_bayesiangssslam}) or semantic
importance---remains an open problem.

\subsection{Dynamic Scene Handling}

Most Gaussian SLAM systems assume static scenes.  WildGS-SLAM addresses
dynamic distractors via uncertainty estimation, and DynActiveGS decomposes
uncertainty into structural vs.\ motion-induced fields.  However, a
complete dynamic-scene Gaussian SLAM with online moving-object tracking,
map segmentation, and consistent static-scene reconstruction is still
lacking.

\subsection{Edge and Mobile Deployment}

The ultimate goal of multimodal Gaussian SLAM is deployment on
resource-constrained platforms: drones, AR headsets, smartphones.  GS-LIVO's
Jetson deployment is a milestone, but real-time 3DGS on mobile-class GPUs
requires fundamentally lighter representations.  The local-primitive
observation of LoCoSplat~\cite{wang2026_locosplat} (0.14M-param MLP replacing
heavy 3D networks) and the budgeted LOD of Budgeted-GS point toward
practical on-device Gaussian SLAM.

\subsection{Evaluation Standardization}

The field lacks standardized benchmarks for multimodal Gaussian SLAM.
LiDAR-based systems evaluate on MARS-LVIG and KITTI; monocular systems on
Replica, TUM, and ScanNet; iToF systems on ZJUL5.  A unified benchmark with
synchronized multi-sensor data (LiDAR, iToF, RGB, IMU, ground truth) across
indoor/outdoor/dynamic scenes would enable fair cross-category comparison.
Gaussian-LIC2's dataset and RMGS-SLAM's looped outdoor benchmark are steps
in this direction.

\subsection{Bidirectional Map--Odometry Coupling}

Gaussian-LIC2 pioneered feeding photometric constraints from the Gaussian
map back into the odometry, demonstrating improved tracking under LiDAR
degeneration.  This bidirectional coupling---where the map is not merely a
passive visualization but an active constraint source---is a promising but
underexplored direction.  Extending it to visual-inertial systems (where the
Gaussian map could provide depth supervision for the tracking front-end) and
to iToF systems (where the Gaussian map could guide ToF depth completion)
could yield significant robustness gains.
